% !TeX root = ../main_ustc.tex
\chapter{总结与展望}
\label{chap:conclusion}

本文围绕约束操作场景下的人机协同控制问题，研究了合作博弈、动态仲裁、几何约束执行和柔性接触力--位协同之间的关系。全文的主线可以概括为一句话：在人能够提供监督经验、机器人需要稳定执行、环境又不断施加几何和力安全约束的任务中，控制器不应简单地选择“听人”或“听机器人”，也不应固定地选择“重位置”或“重力”，而应根据任务阶段、操作者输入和接触状态进行连续仲裁。

\section{全文工作总结}
\label{sec:ch6_summary}

第一章从工业柔性接触任务和机器人辅助手术训练任务出发，分析了操作者意图、机器人自主目标、固定点--轴线几何约束和接触力安全之间的矛盾。通过相关研究综述可以看到，共享控制、微分博弈、阻抗控制、混合力--位控制和安全自主操作已经分别给出了重要基础，但这些方法之间仍存在接口问题。人的输入需要被解释和筛选，接触力边界需要被实时感知和调节，几何约束还需要在工具到机器人法兰的映射中得到保持。基于这一判断，本文将合作博弈和动态仲裁作为统一方法线索。

第二章梳理了全文所需的预备知识。该章统一了机器人任务空间建模、固定点--轴线几何约束、柔性接触 Kelvin--Voigt 模型、线性二次型最优控制、合作微分博弈和 Lyapunov 稳定性分析等基础内容。特别地，第二章明确了全文核心符号的分工：$\alpha_{\mathrm{HR}}$ 表示参考层的人机仲裁权重，$\alpha_{\mathrm{FP}}$ 表示执行层的力--位优先级，$\rho_F$ 表示接触力安全裕度，$\hat K_e$ 表示在线估计的环境等效刚度。这样的符号分层为后续章节避免概念混淆奠定了基础。

第三章面向刚性几何约束任务，提出了意图驱动的合作微分博弈共享控制方法。该章把操作者输入解释为监督参考，而不是直接作为底层命令；再通过交互强度、持续时间和空间风险等特征调节 $\alpha_{\mathrm{HR}}$，使机器人自主参考和人侧监督参考形成连续权衡。几何约束执行部分通过工具尖端到法兰的映射保持固定点--轴线约束，从而使人类修正能够在满足约束的前提下进入机器人任务。

第四章面向柔性接触扫描任务，提出了折扣帕累托力--位协同控制方法。该章建立了包含位置误差、速度误差、接触力误差和力误差泄漏积分的增广系统，并在位置目标和力目标之间构造帕累托标量化代价。与仅依赖固定力--位权重的方法相比，第四章进一步引入力安全裕度 $\rho_F$、边界方向变量 $s_F$、阶段先验和刚度响应项，使 $\alpha_{\mathrm{FP}}$ 能够随接触力风险和环境刚度变化而调整。仿真结果表明，动态仲裁能够在不牺牲主要轨迹质量的前提下降低力抖动和固定权重失配带来的风险。

第五章进一步把第三章的人机意图仲裁和第四章的力--位安全仲裁融合起来，提出面向工业约束接触任务的双仲裁合作博弈控制方法。该章的关键不是把所有目标堆入一个更大的权重，而是保持两层物理含义清楚：$\alpha_{\mathrm{HR}}$ 只作用于参考轨迹生成层，用于决定人类监督修正如何影响 $x_d^I$；$\alpha_{\mathrm{FP}}$ 只作用于执行控制层，用于决定当前位置和接触力目标如何分配优先级。通过安全投影、切向/法向分解、力安全裕度仲裁和基于 $\hat K_e$ 的增益调度，第五章展示了融合方法在安全切向修正、不安全法向压入和变刚度接触中的综合优势。

综上，本文主要完成了以下工作：
\begin{enumerate}[label=(\arabic*)]
  \item 建立了面向约束接触任务的人机协同控制统一符号体系，将人机仲裁、力--位仲裁、几何约束和环境刚度估计分别赋予明确物理含义。
  \item 提出了刚性几何约束下的意图驱动合作博弈共享控制方法，使操作者输入能够以平滑、可解释的方式影响机器人期望轨迹。
  \item 提出了柔性接触任务中的力安全裕度力--位仲裁方法，使控制器能够根据接触力边界风险和局部环境刚度调节位置/力优先级。
  \item 提出了工业约束接触任务中的双仲裁融合框架，在参考层接受安全的人类切向修正，在执行层抑制危险法向压入并保持接触力安全。
\end{enumerate}

\section{不足与展望}
\label{sec:ch6_future}

本文仍存在一些需要继续完善的地方。首先，当前实验主要围绕可控仿真环境和典型柔性接触任务展开，真实工件中的摩擦、局部脱粘、曲率突变、工具磨损和传感器漂移会带来更复杂的残差项。后续工作需要在真实密封条检测、复合材料打磨或手术训练平台上进行更长时间和更多重复次数的验证。

其次，本文的人机意图识别仍以低维交互特征和模糊规则为主。这种设计可解释、易调试，但对复杂语义任务的表达能力有限。未来可以引入视觉检测、语音指令、工艺区域标注和学习式意图预测，使 $\alpha_{\mathrm{HR}}$ 不只依赖主端输入强度，也能理解操作者为什么要干预。

再次，本文的力安全机制主要通过安全投影、边界势函数、硬阈值监督和动态 $\alpha_{\mathrm{FP}}$ 调度实现。这些设计具有明确工程意义，但与严格控制屏障函数相比，理论上的硬安全保证仍不够充分。后续可将控制屏障函数、鲁棒模型预测控制或安全强化学习与本文的合作博弈反馈律结合，使接触力边界和几何约束获得更严格的可证明安全性。

最后，本文将 $\alpha_{\mathrm{HR}}$ 和 $\alpha_{\mathrm{FP}}$ 分层处理，以降低在线求解复杂度。该结构适合当前百赫兹级控制循环，但在更复杂的多接触、多工具或双臂协作任务中，参考层和执行层之间可能存在更强耦合。未来可以研究低维参数化增益表、神经近似黎卡提解和在线凸优化相结合的方法，在保持实时性的同时提升多目标协同能力。

总体来看，约束接触任务中的人机协同控制不是一个单纯追求更大自主性的问题，而是一个需要持续解释、持续仲裁、持续守住安全边界的问题。本文所提出的合作博弈与双仲裁框架只是其中一步，但它提供了一条相对清晰的路线：让人的经验进入参考，让机器人的稳定性守住执行，让接触力和几何约束在控制器内部拥有明确位置。沿着这条路线继续推进，机器人有望在更复杂、更柔性、更安全敏感的工业和医疗操作场景中发挥可靠作用。
